\documentclass[11pt]{article}
\usepackage{amsmath, amssymb, geometry, amsthm, array, tabularx, booktabs, xcolor}
\geometry{letterpaper, margin=1in}

\title{\textbf{Homework 2 Logic: Statements, Negations, and Connectives}}
\author{Tashfeen Omran}
\date{September 2026}

\begin{document}

\maketitle
\tableofcontents
\newpage

% ---------------------------------------------------------
% PART 1: Comprehensive Introduction, Context, and Prerequisites
% ---------------------------------------------------------
\section{Part 1: Comprehensive Introduction, Context, and Prerequisites}

\subsection{Core Concepts}
At the very heart of mathematics and formal reasoning lies \textbf{Logic}. In sections 2.1 through 2.3, we build the fundamental vocabulary required to read, write, and evaluate mathematical propositions. The core ideas are:
\begin{itemize}
    \item \textbf{Statements:} A declarative sentence or assertion that has a definitive truth value—it is either identically true ($T$) or identically false ($F$), but never both. 
    \item \textbf{Open Sentences:} A declarative sentence containing one or more variables. Its truth value depends on the specific values substituted into the variables from a defined \textit{domain}. Once a value is substituted, the open sentence becomes a statement.
    \item \textbf{Negations ($\sim P$):} The logical opposite of a statement. If a statement $P$ is true, its negation $\sim P$ is false, and vice versa. 
    \item \textbf{Conjunctions ($P \land Q$):} The logical ``AND''. A compound statement that is true if and only if \textit{both} $P$ and $Q$ are true.
    \item \textbf{Disjunctions ($P \lor Q$):} The logical (inclusive) ``OR''. A compound statement that is true if at least one of $P$ or $Q$ is true. 
\end{itemize}

\subsection{Intuition and Derivation}
Think of logic as a digital circuit or a binary system. A standard sentence in English can be ambiguous, rhetorical, or subjective (e.g., ``This is a beautiful day.''). Mathematical logic strips away subjectivity. We evaluate sentences purely on their binary state: True ($1$) or False ($0$). 
When we combine statements using connectives like $\land$ (AND) and $\lor$ (OR), we are effectively mapping out intersecting realities. The truth table acts as the ``coordinate system'' for these realities, explicitly defining the outcome of every possible binary combination. The logical ``OR'' is inclusive (meaning $A$, $B$, or both), unlike the casual conversational ``or'' which often implies exclusivity (XOR).

\subsection{Historical Context and Motivation}
The formalization of logic traces back to Aristotle in ancient Greece, who first categorized forms of reasoning and syllogisms. However, the symbolic logic we study today—where sentences are replaced by variables ($P$, $Q$) and operators ($\land$, $\lor$, $\sim$)—was deeply advanced in the 19th century by mathematician George Boole. Boole's motivation was to translate the mechanics of human thought and philosophical deduction into a rigorous algebraic framework (now known as Boolean Algebra). This was driven by the desire to eliminate the ambiguities of natural language in mathematical proofs. Today, this bedrock not only serves pure mathematics but forms the entire theoretical foundation of computer science and digital hardware design.

\subsection{Key Formulas and Truth Tables}
While there are no traditional algebraic ``formulas'' in this section, the foundational ``equations'' of logic are truth tables.

\textbf{1. Negation ($\sim P$)}
\begin{center}
\begin{tabular}{|c|c|}
\hline
$P$ & $\sim P$ \\
\hline
$T$ & $F$ \\
$F$ & $T$ \\
\hline
\end{tabular}
\end{center}

\textbf{2. Conjunction ($P \land Q$) and Disjunction ($P \lor Q$)}
\begin{center}
\begin{tabular}{|c|c|c|c|}
\hline
$P$ & $Q$ & $P \land Q$ & $P \lor Q$ \\
\hline
$T$ & $T$ & $T$ & $T$ \\
$T$ & $F$ & $F$ & $T$ \\
$F$ & $T$ & $F$ & $T$ \\
$F$ & $F$ & $F$ & $F$ \\
\hline
\end{tabular}
\end{center}

\subsection{Prerequisites}
To succeed in this chapter, you must be comfortable with:
\begin{itemize}
    \item \textbf{Basic Set Theory:} Understanding elements ($\in$), subsets ($\subseteq$, $\subset$), unions ($\cup$), intersections ($\cap$), and the empty set ($\emptyset$).
    \item \textbf{Number Classifications:} Knowing the definitions of integers ($\mathbb{Z}$), natural numbers ($\mathbb{N}$), rational numbers ($\mathbb{Q}$), real numbers ($\mathbb{R}$), primes, evens, and odds.
    \item \textbf{Basic Algebra:} Solving simple linear and quadratic equations (e.g., finding the roots of a polynomial).
\end{itemize}

\newpage
% ---------------------------------------------------------
% PART 2: Detailed Homework Solutions
% ---------------------------------------------------------
\section{Part 2: Detailed Homework Solutions}

\subsection{Section 2.1 Exercises}

\textbf{2.1. Which of the following sentences are statements? For those that are, indicate the truth value.}
\begin{itemize}
    \item[(a)] The integer $123$ is prime. \\
    \textbf{Solution:} This is a declarative sentence with a definitive truth value. It is a \textbf{statement}. Because $123 = 3 \times 41$, it is not prime. Its truth value is \textbf{False ($F$)}.
    
    \item[(b)] The integer $0$ is even. \\
    \textbf{Solution:} This is a declarative sentence. It is a \textbf{statement}. Since $0 = 2 \times 0$ (an integer multiple of $2$), it is even. Its truth value is \textbf{True ($T$)}.
    
    \item[(c)] Is $5 \times 2 = 10$? \\
    \textbf{Solution:} This is an interrogative sentence (a question). It is \textbf{not a statement}.
    
    \item[(d)] $x^2 - 4 = 0$. \\
    \textbf{Solution:} This sentence contains a variable $x$ and no domain or specific value is assigned. Therefore, it is an open sentence, \textbf{not a statement}.
    
    \item[(e)] Multiply $5x + 2$ by $3$. \\
    \textbf{Solution:} This is an imperative sentence (a command). It is \textbf{not a statement}.
    
    \item[(f)] $5x + 3$ is an odd integer. \\
    \textbf{Solution:} This contains a variable $x$ whose value is unspecified. Thus, it is an open sentence, \textbf{not a statement}.
    
    \item[(g)] What an impossible question! \\
    \textbf{Solution:} This is an exclamatory sentence. It is \textbf{not a statement}.
\end{itemize}

\vspace{0.5cm}
\textbf{2.2. Consider the sets $A, B, C$ and $D$ below. Which of the following statements are true? Give an explanation for each false statement.}
\[ A = \{1, 4, 7, 10, 13, 16, \dots\} \quad C = \{x \in \mathbb{Z} : x \text{ is prime and } x \neq 2\} \]
\[ B = \{x \in \mathbb{Z} : x \text{ is odd}\} \quad D = \{1, 2, 3, 5, 8, 13, 21, 34, 55, \dots\} \]

\begin{itemize}
    \item[(a)] $25 \in A$ \\
    \textbf{Solution:} The set $A$ consists of positive integers of the form $3k + 1$ for integers $k \ge 0$. Notice that $25 = 3(8) + 1$. Thus, $25 \in A$. This statement is \textbf{True}.
    
    \item[(b)] $33 \in D$ \\
    \textbf{Solution:} The set $D$ is the sequence of Fibonacci numbers. The numbers in the sequence around $33$ are $21, 34$, and $55$. Since $33$ is not one of them, the statement is \textbf{False}. Explanation: $33$ is not a Fibonacci number.
    
    \item[(c)] $22 \notin A \cup D$ \\
    \textbf{Solution:} First, check if $22 \in A$. Notice $22 = 3(7) + 1$, so $22 \in A$. Because $22$ is in $A$, it is in the union $A \cup D$. Therefore, the statement $22 \notin A \cup D$ is \textbf{False}. Explanation: $22$ is an element of $A$, so it must be an element of $A \cup D$.
    
    \item[(d)] $C \subseteq B$ \\
    \textbf{Solution:} The set $C$ contains all prime numbers except $2$. Since $2$ is the only even prime number, all elements in $C$ are odd primes. Therefore, every element in $C$ is an odd integer, meaning every element is in $B$. This statement is \textbf{True}.
    
    \item[(e)] $\emptyset \in B \cap D$ \\
    \textbf{Solution:} The intersection $B \cap D$ contains integers that are both odd and Fibonacci numbers (e.g., $1, 3, 5, 13, 21$). The empty set $\emptyset$ is a set, not an integer. Therefore, $\emptyset$ cannot be an element of $B \cap D$. The statement is \textbf{False}. Explanation: The empty set is a set, not an integer, so it cannot be a member of a set of integers.
    
    \item[(f)] $53 \notin C$ \\
    \textbf{Solution:} The integer $53$ is a prime number and $53 \neq 2$. Therefore, $53 \in C$. Thus, the claim that $53 \notin C$ is \textbf{False}. Explanation: $53$ is an odd prime number, meaning it satisfies the condition to be in set $C$.
\end{itemize}

\vspace{0.5cm}
\textbf{2.3. Which of the following statements are true? Give an explanation for each false statement.}
\begin{itemize}
    \item[(a)] $\emptyset \in \emptyset$ \\
    \textbf{Solution:} The empty set, by definition, contains absolutely no elements. Therefore, it cannot contain anything, including the empty set itself. This is \textbf{False}. Explanation: The empty set has no elements.
    
    \item[(b)] $\emptyset \in \{\emptyset\}$ \\
    \textbf{Solution:} The set $\{\emptyset\}$ is a set containing exactly one element: the empty set. Therefore, $\emptyset$ is an element of this set. This is \textbf{True}.
    
    \item[(c)] $\{1, 3\} = \{3, 1\}$ \\
    \textbf{Solution:} In set theory, the order of elements does not matter. Both sets contain exactly the same distinct elements. This is \textbf{True}.
    
    \item[(d)] $\emptyset = \{\emptyset\}$ \\
    \textbf{Solution:} The set on the left has $0$ elements. The set on the right has $1$ element (which happens to be the empty set). Two sets of different sizes cannot be equal. This is \textbf{False}. Explanation: $\emptyset$ contains no elements, whereas $\{\emptyset\}$ contains one element.
    
    \item[(e)] $\emptyset \subset \{\emptyset\}$ \\
    \textbf{Solution:} The empty set is a proper subset of any non-empty set. Since $\{\emptyset\}$ is non-empty, $\emptyset$ is a proper subset of it. This is \textbf{True}.
    
    \item[(f)] $1 \subseteq \{1\}$ \\
    \textbf{Solution:} The symbol $\subseteq$ denotes a subset relationship, which must occur between two sets. The integer $1$ is an element, not a set. (The correct statement would be $1 \in \{1\}$ or $\{1\} \subseteq \{1\}$). This is \textbf{False}. Explanation: $1$ is an element, not a set, so it cannot be a subset.
\end{itemize}

\vspace{0.5cm}
\textbf{2.4. Consider the open sentence $P(x) : x(x - 1) = 6$ over the domain $\mathbb{R}$.}
\begin{itemize}
    \item[(a)] For what values of $x$ is $P(x)$ a true statement? \\
    \textbf{Solution:} We solve the algebraic equation:
    \begin{align*}
        x(x - 1) &= 6 \\
        x^2 - x - 6 &= 0 \\
        (x - 3)(x + 2) &= 0
    \end{align*}
    Thus, $x = 3$ or $x = -2$. $P(x)$ is a true statement for \textbf{$x \in \{3, -2\}$}.
    
    \item[(b)] For what values of $x$ is $P(x)$ a false statement? \\
    \textbf{Solution:} $P(x)$ is false for any real number that is not a solution to the equation. Thus, it is false for all \textbf{$x \in \mathbb{R} \setminus \{3, -2\}$}.
\end{itemize}

\vspace{0.5cm}
\textbf{2.5. For the open sentence $P(x): 3x - 2 > 4$ over the domain $\mathbb{Z}$, determine:}
\begin{itemize}
    \item[(a)] the values of $x$ for which $P(x)$ is true. \\
    \textbf{Solution:} Solve the inequality:
    \begin{align*}
        3x - 2 &> 4 \\
        3x &> 6 \\
        x &> 2
    \end{align*}
    Since the domain is $\mathbb{Z}$, $P(x)$ is true for integers strictly greater than 2: \textbf{$x \in \{3, 4, 5, 6, \dots\}$}.
    
    \item[(b)] the values of $x$ for which $P(x)$ is false. \\
    \textbf{Solution:} $P(x)$ is false for all integers less than or equal to 2: \textbf{$x \in \{\dots, -2, -1, 0, 1, 2\}$}.
\end{itemize}

\vspace{0.5cm}
\textbf{2.6. For the open sentence $Q(A): A \subseteq \{1, 2, 3\}$ over the domain $S = \mathcal{P}(\{1, 2, 4\})$, determine:} \\
First, list the elements of the domain $S$ (the power set of $\{1, 2, 4\}$): 
\[ S = \{\emptyset, \{1\}, \{2\}, \{4\}, \{1, 2\}, \{1, 4\}, \{2, 4\}, \{1, 2, 4\}\} \]

\begin{itemize}
    \item[(a)] all $A \in S$ for which $Q(A)$ is true. \\
    \textbf{Solution:} We need subsets of $\{1, 2, 4\}$ that are also subsets of $\{1, 2, 3\}$. This means $A$ can only contain the elements $1$ and $2$. These are: \textbf{$\emptyset, \{1\}, \{2\}, \{1, 2\}$}.
    
    \item[(b)] all $A \in S$ for which $Q(A)$ is false. \\
    \textbf{Solution:} This applies to any set in $S$ that contains the element $4$, as $4$ is not in $\{1, 2, 3\}$. These are: \textbf{$\{4\}, \{1, 4\}, \{2, 4\}, \{1, 2, 4\}$}.
    
    \item[(c)] all $A \in S$ for which $A \cap \{1, 2, 3\} = \emptyset$. \\
    \textbf{Solution:} This requires $A$ to share no elements with $\{1, 2, 3\}$. The only allowed element from $S$ is $4$. The sets in $S$ composed entirely of elements not in $\{1, 2, 3\}$ are: \textbf{$\emptyset, \{4\}$}.
\end{itemize}

\vspace{0.5cm}
\textbf{2.7. Let $P(n) : n$ and $n+2$ are primes, be an open sentence over the domain $\mathbb{N}$. Find six positive integers $n$ for which $P(n)$ is true. If $n \in \mathbb{N}$ such that $P(n)$ is true, then the two integers $n, n+2$ are called twin primes.} \\
\textbf{Solution:} We must find pairs of prime numbers that differ by exactly 2. Testing small integers:
\begin{itemize}
    \item $n=3$: $3$ and $5$ are prime (True).
    \item $n=5$: $5$ and $7$ are prime (True).
    \item $n=11$: $11$ and $13$ are prime (True).
    \item $n=17$: $17$ and $19$ are prime (True).
    \item $n=29$: $29$ and $31$ are prime (True).
    \item $n=41$: $41$ and $43$ are prime (True).
\end{itemize}
Thus, six integers are \textbf{$n \in \{3, 5, 11, 17, 29, 41\}$}.

\vspace{0.5cm}
\textbf{2.8. Let $P(n): \frac{n^2+5n+6}{2}$ is even.}
\begin{itemize}
    \item[(a)] Find a set $S_1$ of three integers such that $P(n)$ is an open sentence over the domain $S_1$ and $P(n)$ is true for each $n \in S_1$. \\
    \textbf{Solution:} Simplify the expression: $\frac{n^2+5n+6}{2} = \frac{(n+2)(n+3)}{2}$. Let's test a few integers.
    \begin{align*}
        n=1 &\implies \frac{(3)(4)}{2} = 6 \quad (\text{even, True}) \\
        n=2 &\implies \frac{(4)(5)}{2} = 10 \quad (\text{even, True}) \\
        n=3 &\implies \frac{(5)(6)}{2} = 15 \quad (\text{odd, False}) \\
        n=4 &\implies \frac{(6)(7)}{2} = 21 \quad (\text{odd, False}) \\
        n=5 &\implies \frac{(7)(8)}{2} = 28 \quad (\text{even, True})
    \end{align*}
    Based on our tests, we can choose \textbf{$S_1 = \{1, 2, 5\}$}.
    
    \item[(b)] Find a set $S_2$ of three integers such that $P(n)$ is an open sentence over the domain $S_2$ and $P(n)$ is false for each $n \in S_2$. \\
    \textbf{Solution:} Continuing the pattern, let's test $n=6, 7$.
    \begin{align*}
        n=6 &\implies \frac{(8)(9)}{2} = 36 \quad (\text{even, True}) \\
        n=7 &\implies \frac{(9)(10)}{2} = 45 \quad (\text{odd, False})
    \end{align*}
    From the tests where the result is odd, we can choose \textbf{$S_2 = \{3, 4, 7\}$}.
\end{itemize}

\vspace{0.5cm}
\textbf{2.9. Find an open sentence $P(n)$ over the domain $S = \{3, 5, 7, 9\}$ such that $P(n)$ is true for half of the integers in $S$ and false for the other half.} \\
\textbf{Solution:} We need an open sentence that holds true for exactly two elements. For example, $n$ being less than or equal to $6$ is true for $3$ and $5$, but false for $7$ and $9$. 
\textbf{Answer: $P(n): n \le 6$.} (Many other answers are possible, e.g., $P(n): n \text{ is a prime number.}$)

\vspace{0.5cm}
\textbf{2.10. Find two open sentences $P(n)$ and $Q(n)$, both over the domain $S = \{2, 4, 6, 8\}$, such that $P(2)$ and $Q(2)$ are both true, $P(4)$ and $Q(4)$ are both false, $P(6)$ is true and $Q(6)$ is false, while $P(8)$ is false and $Q(8)$ is true.} \\
\textbf{Solution:} We summarize the requirements:
\begin{center}
\begin{tabular}{ccc}
\hline
$n$ & $P(n)$ & $Q(n)$ \\
\hline
2 & True & True \\
4 & False & False \\
6 & True & False \\
8 & False & True \\
\hline
\end{tabular}
\end{center}
Notice $P(n)$ is true for $\{2, 6\}$. These are numbers in $S$ that are not divisible by $4$. Thus, $P(n): \frac{n}{2}$ is an odd integer.
Notice $Q(n)$ is true for $\{2, 8\}$. Thus, we can simply define $Q(n)$ by listing the valid cases: $Q(n): n = 2 \text{ or } n = 8$.
\textbf{Answer: $P(n): \frac{n}{2} \text{ is an odd integer.}$ and $Q(n): n=2 \text{ or } n=8$.}

\subsection{Section 2.2 Exercises}

\textbf{2.11. State the negation of each of the following statements.}
\begin{itemize}
    \item[(a)] $\sqrt{2}$ is a rational number. \\
    \textbf{Solution: $\sqrt{2}$ is not a rational number.} (Or equivalently, $\sqrt{2}$ is an irrational number).
    
    \item[(b)] $0$ is not a negative integer. \\
    \textbf{Solution: $0$ is a negative integer.}
    
    \item[(c)] $111$ is a prime number. \\
    \textbf{Solution: $111$ is not a prime number.}
\end{itemize}

\vspace{0.5cm}
\textbf{2.12. Complete the truth table in Figure 2.3.} \\
\textbf{Solution:}
\begin{center}
\begin{tabular}{|c|c|c|c|}
\hline
$P$ & $Q$ & $\sim P$ & $\sim Q$ \\
\hline
$T$ & $T$ & $F$ & $F$ \\
$T$ & $F$ & $F$ & $T$ \\
$F$ & $T$ & $T$ & $F$ \\
$F$ & $F$ & $T$ & $T$ \\
\hline
\end{tabular}
\end{center}

\vspace{0.5cm}
\textbf{2.13. State the negation of each of the following statements.}
\begin{itemize}
    \item[(a)] The real number $r$ is at most $\sqrt{2}$. \\
    \textbf{Solution: The real number $r$ is greater than $\sqrt{2}$.}
    
    \item[(b)] The absolute value of the real number $a$ is less than $3$. \\
    \textbf{Solution: The absolute value of the real number $a$ is at least $3$.}
    
    \item[(c)] Two angles of the triangle are $45^\circ$. \\
    \textbf{Solution: It is not the case that two angles of the triangle are $45^\circ$.} (Or, ``The triangle does not have two angles measuring $45^\circ$.")
    
    \item[(d)] The area of the circle is at least $9\pi$. \\
    \textbf{Solution: The area of the circle is less than $9\pi$.}
    
    \item[(e)] Two sides of the triangle have the same length. \\
    \textbf{Solution: No two sides of the triangle have the same length.}
    
    \item[(f)] The point $P$ in the plane lies outside of the circle $C$. \\
    \textbf{Solution: The point $P$ in the plane lies on or inside the circle $C$.}
\end{itemize}

\vspace{0.5cm}
\textbf{2.14. State the negation of each of the following statements.}
\begin{itemize}
    \item[(a)] At least two of my library books are overdue. \\
    \textbf{Solution: At most one of my library books is overdue.}
    
    \item[(b)] One of my two friends misplaced his homework assignment. \\
    \textbf{Solution: Neither of my two friends misplaced his homework assignment.} (Assuming the original implied "at least one").
    
    \item[(c)] No one expected that to happen. \\
    \textbf{Solution: At least one person expected that to happen.}
    
    \item[(d)] It's not often that my instructor teaches that course. \\
    \textbf{Solution: It is often that my instructor teaches that course.}
    
    \item[(e)] It's surprising that two students received the same exam score. \\
    \textbf{Solution: It is not surprising that two students received the same exam score.}
\end{itemize}

\subsection{Section 2.3 Exercises}

\textbf{2.15. Complete the truth table in Figure 2.6.} \\
\textbf{Solution:}
\begin{center}
\begin{tabular}{|c|c|c|c|}
\hline
$P$ & $Q$ & $\sim Q$ & $P \land (\sim Q)$ \\
\hline
$T$ & $T$ & $F$ & $F$ \\
$T$ & $F$ & $T$ & $T$ \\
$F$ & $T$ & $F$ & $F$ \\
$F$ & $F$ & $T$ & $F$ \\
\hline
\end{tabular}
\end{center}

\vspace{0.5cm}
\textbf{2.16. For the sets $A = \{1, 2, \dots, 10\}$ and $B = \{2, 4, 6, 9, 12, 25\}$, consider the statements: $P: A \subseteq B \quad Q: |A \setminus B| = 6$. Determine which of the following statements are true.} \\
First, evaluate the truth values of $P$ and $Q$.
$A$ contains $1$, which is not in $B$. Thus, $A \subseteq B$ is \textbf{False} ($P$ is $F$).
The set difference $A \setminus B$ (elements in $A$ not in $B$) is $\{1, 3, 5, 7, 8, 10\}$. There are $6$ elements here. Thus, $|A \setminus B| = 6$ is \textbf{True} ($Q$ is $T$).
\begin{itemize}
    \item[(a)] $P \lor Q \implies F \lor T = \textbf{True}$
    \item[(b)] $P \lor (\sim Q) \implies F \lor F = \textbf{False}$
    \item[(c)] $P \land Q \implies F \land T = \textbf{False}$
    \item[(d)] $(\sim P) \land Q \implies T \land T = \textbf{True}$
    \item[(e)] $(\sim P) \lor (\sim Q) \implies T \lor F = \textbf{True}$
\end{itemize}

\vspace{0.5cm}
\textbf{2.17. Let $P: 15$ is odd, and $Q: 21$ is prime. State each of the following in words and determine whether it is true or false.} \\
First, evaluate $P$ and $Q$. $15$ is odd, so $P$ is \textbf{True}. $21$ is $3 \times 7$, so $21$ is not prime. $Q$ is \textbf{False}.
\begin{itemize}
    \item[(a)] $P \lor Q$ \\
    \textbf{Words:} 15 is odd or 21 is prime. \\
    \textbf{Truth Value:} $T \lor F = \textbf{True}$.
    
    \item[(b)] $P \land Q$ \\
    \textbf{Words:} 15 is odd and 21 is prime. \\
    \textbf{Truth Value:} $T \land F = \textbf{False}$.
    
    \item[(c)] $(\sim P) \lor Q$ \\
    \textbf{Words:} 15 is not odd or 21 is prime. \\
    \textbf{Truth Value:} $F \lor F = \textbf{False}$.
    
    \item[(d)] $P \land (\sim Q)$ \\
    \textbf{Words:} 15 is odd and 21 is not prime. \\
    \textbf{Truth Value:} $T \land T = \textbf{True}$.
\end{itemize}

\vspace{0.5cm}
\textbf{2.18. Let $S = \{1, 2, \dots, 6\}$ and let $P(A): A \cap \{2, 4, 6\} = \emptyset$, and $Q(A): A \neq \emptyset$, be open sentences over the domain $\mathcal{P}(S)$.} \\
Note: $\mathcal{P}(S)$ is the power set (all subsets) of $S$. $P(A)$ means subset $A$ contains no even numbers. $Q(A)$ means subset $A$ is not empty.
\begin{itemize}
    \item[(a)] Determine all $A \in \mathcal{P}(S)$ for which $P(A) \land Q(A)$ is true. \\
    \textbf{Solution:} $A$ must contain no even numbers (so it must be formed purely from the odd numbers $\{1, 3, 5\}$) AND it must not be empty. 
    \textbf{Answer:} $A \in \{\{1\}, \{3\}, \{5\}, \{1, 3\}, \{1, 5\}, \{3, 5\}, \{1, 3, 5\}\}$.
    
    \item[(b)] Determine all $A \in \mathcal{P}(S)$ for which $P(A) \lor (\sim Q(A))$ is true. \\
    \textbf{Solution:} $A$ must contain no even numbers OR $A$ must be empty. Since the empty set already contains no even numbers, this condition covers all subsets made entirely of $\{1, 3, 5\}$, including the empty set.
    \textbf{Answer:} $A \in \{\emptyset, \{1\}, \{3\}, \{5\}, \{1, 3\}, \{1, 5\}, \{3, 5\}, \{1, 3, 5\}\}$.
    
    \item[(c)] Determine all $A \in \mathcal{P}(S)$ for which $(\sim P(A)) \land (\sim Q(A))$ is true. \\
    \textbf{Solution:} $A$ must contain at least one even number AND $A$ must be empty. This is a logical impossibility. 
    \textbf{Answer: None (or $\emptyset$).}
\end{itemize}

\newpage
% ---------------------------------------------------------
% PART 3: In-Depth Analysis of Problems and Techniques
% ---------------------------------------------------------
\section{Part 3: In-Depth Analysis of Problems and Techniques}

\subsection{A) Problem Types and General Approach}
\begin{enumerate}
    \item \textbf{Statement Identification (Ex 2.1):} 
    \textit{Approach:} Scan the sentence for absolute declaration. If it asks a question, gives a command, or uses an unspecified variable without a domain, it's immediately disqualified as a mathematical statement. 
    \item \textbf{Set Membership and Inclusion (Ex 2.2, 2.3):}
    \textit{Approach:} Distinguish fiercely between an element ($\in$) and a subset ($\subseteq$). A subset is a collection of elements, whereas an element is the object itself. 
    \item \textbf{Evaluating Open Sentences (Ex 2.4, 2.5, 2.6, 2.8):}
    \textit{Approach:} Solve the given algebraic constraint ($x^2 - x - 6 = 0$, $3x - 2 > 4$). The solution set defines the exact inputs within the domain that render the open sentence \textbf{True}. All other inputs within the domain render it \textbf{False}.
    \item \textbf{English-to-Logic Translation for Negations (Ex 2.11, 2.13, 2.14):}
    \textit{Approach:} Pinpoint the operational inequality or quantifier word. ``At least $X$'' negates to ``Strictly less than $X$''. ``No one'' negates to ``At least one person''. Be careful not to alter the subject, only the truth condition.
\end{enumerate}

\subsection{B) Key Algebraic and Logical Manipulations}
\begin{itemize}
    \item \textbf{Algebraic Factoring to Find Truth Domains:} In Ex 2.4, we encountered $x(x - 1) = 6$. The crucial trick was to move the $6$ over and factor to $x^2 - x - 6 = (x-3)(x+2) = 0$. In logic, finding when a statement is True often perfectly mirrors finding the roots of an equation.
    \item \textbf{Parity Analysis (Odd/Even testing):} In Ex 2.8, determining if $\frac{(n+2)(n+3)}{2}$ is even required checking if the product of the sequential integers contained a factor of $4$. We systematically tested integers $n=1,2,3,4,5$ to establish a visible alternating pattern. 
    \item \textbf{Power Set Intersections:} In Ex 2.18, dealing with $\mathcal{P}(S)$, recognizing that $A \cap \{2,4,6\} = \emptyset$ is equivalent to saying ``$A$ is constructed strictly from the complement of $\{2,4,6\}$ in $S$'' reduced a massive $64$-element space search down to just the $8$ subsets of $\{1,3,5\}$.
\end{itemize}

% ---------------------------------------------------------
% PART 4: "Cheatsheet" and Tips for Success
% ---------------------------------------------------------
\section{Part 4: "Cheatsheet" and Tips for Success}

\subsection{Summary of Formulas}
\begin{itemize}
    \item \textbf{Conjunction ($P \land Q$):} True \textit{only} if both $P$ and $Q$ are True. (Analogy: Set Intersection $\cap$).
    \item \textbf{Disjunction ($P \lor Q$):} False \textit{only} if both $P$ and $Q$ are False. (Analogy: Set Union $\cup$).
    \item \textbf{Negation ($\sim P$):} Flips the truth value.
\end{itemize}

\subsection{Cheats, Tricks, and Shortcuts}
\begin{itemize}
    \item \textbf{The Empty Set Cheat:} The empty set ($\emptyset$) is always a subset of every set. However, it is \textit{rarely} an element of a set unless you are dealing with power sets ($\mathcal{P}(S)$) or sets containing sets (like $\{\emptyset\}$). 
    \item \textbf{Negating Inequalities Shortcut:} When negating quantitative inequalities, flip the sign and reverse the equals inclusion. 
        \begin{itemize}
            \item $>$ becomes $\le$
            \item $<$ becomes $\ge$
            \item $\ge$ becomes $<$
        \end{itemize}
\end{itemize}

\subsection{Common Pitfalls and Traps}
\begin{itemize}
    \item \textbf{Assuming ``Or'' is Exclusive:} In everyday English, ``You can have chicken or fish'' means you can't have both. In mathematics, $P \lor Q$ is \textbf{inclusive}. It is True if $P$ is true, if $Q$ is true, \textit{and if both are true}. 
    \item \textbf{Negating ``At least two'':} Students often mistakenly negate ``At least two'' as ``None''. The correct negation of ``At least two items'' is ``At most one item'' (meaning one or zero).
    \item \textbf{Confusing Open Sentences with Statements:} Just because $x^2 + 1 > 0$ looks like an absolute truth, it is technically an \textit{open sentence} over $\mathbb{R}$. A statement must have its domain and quantifiers explicitly declared or be dealing with constants.
\end{itemize}

% ---------------------------------------------------------
% PART 5: Conceptual Synthesis and The "Big Picture"
% ---------------------------------------------------------
\section{Part 5: Conceptual Synthesis and The "Big Picture"}

\subsection{A) Thematic Connections}
The core theme of this chapter is the \textbf{discretization of truth}. Much like how calculus allows us to zoom in on an infinitely complex curve until it becomes a straight, measurable line, logic forces abstract, ambiguous human sentences into rigid, measurable binary states (True/False). This connects to set theory, where an element is entirely inside a set or completely outside it—there is no middle ground (the Law of Excluded Middle). 

\subsection{B) Forward and Backward Links}
\textbf{Backward:} This material is a formalization of the sets and algebra you learned previously. Solving $3x - 2 > 4$ is basic algebra; defining the solution set as the precise truth domain of an open sentence is the logical upgrade. 
\textbf{Forward:} These basic operators ($\sim, \lor, \land$) act as the bedrock for logical Implications ($P \implies Q$), which are the core of mathematical proofs. You cannot prove ``If a function is differentiable, then it is continuous'' without structurally understanding how ``if-then'' statements are built from conjunctions and negations.

% ---------------------------------------------------------
% PART 6: Real-World Application and Modeling
% ---------------------------------------------------------
\section{Part 6: Real-World Application and Modeling (Finance Focus)}

\subsection{A) Concrete Scenarios}
\begin{enumerate}
    \item \textbf{Algorithmic Trading (Execution Logic):} A quantitative trading program must execute a block trade only if several volatile conditions align perfectly. An execution statement $E$ might be defined by a massive conjunction of sub-statements: $E = P_1 \land P_2 \land (\sim P_3)$, where $P_1$ is ``The 50-day moving average crosses the 200-day'', $P_2$ is ``RSI is below 30'', and $P_3$ is ``The VIX is above 40''. Logic maps this perfectly.
    \item \textbf{Derivative Pricing (Knock-Out Options):} A knock-out option becomes worthless if the underlying asset's price crosses a specific barrier. The payoff condition heavily utilizes set intersections and logical negations. 
    \item \textbf{Credit Default Swaps (CDS):} In structured finance, a portfolio CDS might only trigger a payout if multiple corporate entities default. The trigger condition is a logical disjunction (e.g., $D_A \lor D_B \lor D_C$), and pricing this security requires calculating the joint probability of these overlapping truth domains. 
\end{enumerate}

\subsection{B) Model Problem Setup}
\textbf{Scenario: Algorithmic Hedge Activation} \\
Let $S_T$ represent the price of a market index at time $T$. Let $V_T$ represent the market volatility index (VIX) at time $T$. 
A hedge fund algorithm contains the following open sentences over the domain of time $T > 0$:
\[ P(T): S_T < 4000 \]
\[ Q(T): V_T \ge 30 \]
The fund's risk policy states: \textit{``Activate the downside hedge ($H$) if and only if the index falls below 4000 AND volatility is at least 30.'' }
\textbf{The Logical Model:}
The condition for the hedge to trigger is $H(T) = P(T) \land Q(T)$. 
If we later want to audit when the hedge \textit{failed} to trigger, we must evaluate the negation:
\[ \sim H(T) = \sim(P(T) \land Q(T)) \]
(As you will learn in De Morgan's Laws, this evaluates to the disjunction $\sim P(T) \lor \sim Q(T)$, meaning the hedge stays off as long as the index is at least 4000 OR the volatility is under 30).

% ---------------------------------------------------------
% PART 7: Common Variations and Untested Concepts
% ---------------------------------------------------------
\section{Part 7: Common Variations and Untested Concepts}

The homework set provided covered the basics of $\land$, $\lor$, and $\sim$. However, in standard logic curricula, these concepts are immediately followed by two major logical structures that were absent in these specific exercises:

\textbf{1. Implications ($P \implies Q$)} 
An implication is an "If... then" statement. It is only false when the hypothesis $P$ is true, but the conclusion $Q$ is false. 
\textit{Example:} Let $P$ be ``It is raining'' and $Q$ be ``The ground is wet''. The statement $P \implies Q$ translates to ``If it is raining, then the ground is wet.''
A critical variation is the logical equivalence: $(P \implies Q) \equiv (\sim P) \lor Q$. 

\textbf{2. De Morgan's Laws}
De Morgan's Laws provide the exact algebraic rules for distributing a negation $\sim$ across conjunctions and disjunctions.
\begin{align*}
    \sim (P \land Q) &\equiv (\sim P) \lor (\sim Q) \\
    \sim (P \lor Q) &\equiv (\sim P) \land (\sim Q)
\end{align*}
\textit{Example Problem:} Find the negation of ``I am rich and I am happy.''
\textit{Solution:} By De Morgan's Law, the negation of $(P \land Q)$ is $(\sim P) \lor (\sim Q)$. Thus, the negation is ``I am not rich OR I am not happy.''

% ---------------------------------------------------------
% PART 8: Advanced Diagnostic Testing: "Find the Flaw"
% ---------------------------------------------------------
\section{Part 8: Advanced Diagnostic Testing: "Find the Flaw"}

Below are five logic problems with seemingly correct solutions that contain a subtle, critical error. Your task is to act as the grader, locate the error, explain it in one sentence, and provide the corrected final solution.

\subsection{Problem 1: Statement Identification}
\textbf{Prompt:} Determine if the following sentence is a statement: ``$x^2 + 1 > 0$''. \\
\textbf{Flawed Solution:} Because $x^2$ is always non-negative for any real number, $x^2 + 1$ is always strictly greater than $0$. Since this is identically true forever, it has a definitive truth value of True. Therefore, it is a statement. \\
\textbf{Find the flaw and correct it.}

\subsection{Problem 2: Negating an Inequality}
\textbf{Prompt:} State the negation of the following sentence: ``The temperature $T$ is less than $100$ degrees.'' \\
\textbf{Flawed Solution:} The original statement is $T < 100$. The logical opposite of being less than $100$ is being greater than $100$. Therefore, the negation is $T > 100$, or in words, ``The temperature $T$ is greater than $100$ degrees.'' \\
\textbf{Find the flaw and correct it.}

\subsection{Problem 3: The Disjunction Truth Table}
\textbf{Prompt:} Let $P$ be ``$2+2=4$'' and $Q$ be ``Paris is in France.'' Evaluate the truth value of $P \lor Q$. \\
\textbf{Flawed Solution:} We know that $P$ is True and $Q$ is True. The $\lor$ operator means ``or''. Because both statements are True simultaneously, the ``or'' condition fails (it cannot be one or the other if it is both). Therefore, $P \lor Q$ is False. \\
\textbf{Find the flaw and correct it.}

\subsection{Problem 4: Subset Logic}
\textbf{Prompt:} Let $S = \{1, 2, 3\}$. Let $P(A)$ be the open sentence $A \subseteq \{1\}$ over the domain of the power set of $S$. Is $P(\{1, 2\})$ True or False? \\
\textbf{Flawed Solution:} To check if $\{1, 2\} \subseteq \{1\}$, we check the elements. The set $\{1, 2\}$ contains the element $1$, which perfectly matches the element in $\{1\}$. Since there is an overlapping element, the subset relation holds. Thus, $P(\{1, 2\})$ is True. \\
\textbf{Find the flaw and correct it.}

\subsection{Problem 5: Distributing Negations (De Morgan Error)}
\textbf{Prompt:} Write the negation of the statement: ``The asset's price is rising AND the volume is high.'' \\
\textbf{Flawed Solution:} Let $P$ be ``price is rising'' and $Q$ be ``volume is high''. The statement is $P \land Q$. To negate it, we negate both pieces. The negation is $\sim P \land \sim Q$. In words: ``The asset's price is not rising AND the volume is not high.'' \\
\textbf{Find the flaw and correct it.}

\end{document}